
In labor markets, employers recognize that signals about their preferences will affect the set of sellers that could be potential partners in trade.
For example, the most basic form of signaling is a job posting:
a firm looking for a software engineer will likely attract applicants with software engineering experience.
This form of signaling is related to {\em horizontal differentiation} in the preferences of buyers:
e.g., some employers want software engineers, and others want accountants.
Conveying preferences relating to horizontal differentation is straightforward and lacks much strategic consideration:
a firm looking to hire a software engineer has no incentive to say they are looking for an accountant.
The would-be job applicants and the employer essentially have aligned incentives.

By contrast, conveying preferences related to {\em vertical} differentiation---differences in quality---gives rise to complex incentives.
Nominally, if an employer declares they are willing to pay for the highest skilled talent, reasoning similar to the previous paragraph suggests they would attract only high quality workers to their job openings.
Of course, there is a serious countervailing force:
if an employer has stated that they have a high willingness to pay, but are unable to perfectly ascertain quality, they might find themselves looking at applicants of the same quality as before, but now at a higher price.


This pooling equilibrium is unwelcome because it likely leads to worse matches: 
employers have a more heterogeneous pool of would-be applicants than they would like. 
Giving them a pool of applicants that they would prefer more is socially attractive, even if it also means employers may end up with a relatively worse deal in negotiations (conditional on matching).
From the perspective of the market, efficiency is driven by increasing the quality and quantity of matches.
Changes that lead to buyers getting a relatively worse deal when they do match---but only because they have to pay higher wages---might be very attractive if objectively more and better matches were formed.
However, due to the competing incentives described above, and despite the potential welfare benefits, in practice honest information disclosure cannot be easily compelled in the labor market.

In this paper we study whether employers would {\em endogeneously} choose to signal their true price and quality preferences to the other side of the market, despite the potential adverse consequences.
We consider a market intervention in which buyers in an online labor market were forced to make a claim about their relative preference over price and quality. 



\title{Eliciting Buyer Preferences as a Platform Design Problem\footnote{Author contact information, data sets and code are
     currently or will be available at
     \href{http://www.john-joseph-horton.com/}{http://www.john-joseph-horton.com/}.}}

\title{Eliciting and Conveying Preferences in a Matching Market: \\
  Evidence from a Field Experiment\footnote{Author contact information, data sets and code are
     currently or will be available at
     \href{http://www.john-joseph-horton.com/}{http://www.john-joseph-horton.com/}.}}


\title{Eliciting and Conveying Preferences: \\
  Evidence from a Field Experiment\footnote{Author contact information, data sets and code are
     currently or will be available at
     \href{http://www.john-joseph-horton.com/}{http://www.john-joseph-horton.com/}.}}


\subsection{Discussion of match outcome results}

Our results suggest that the separating equilibrium likely improved match quality, at least as measured by our quantity outcomes. 
This raises the more general question: under what conditions can we expect an overall increase in welfare in the separating equilibrium?
To obtain some insight into this question, we return to the simplified framework introduced in Section \ref{sec:discussion1}.  

If we presume that workers' payoffs are only the wages they earn, then since wages are just transfers, overall market welfare in the equilibrium is simply the average value of $v_t \theta$ over all employers (where $t$ is the employer type and $\theta$ is the quality of the hired worker).  
Suppose there is a unit mass of employers in the market, with a fraction $\mu$ being high type, and $1-\mu$ being low type.  
It follows that aggregate market welfare in the separating equilibrium is:
\begin{align}
\mu v_H q_H + (1 - \mu) v_L q_L. 
\end{align}
On the other hand, suppose in the pooling equilibrium that average hired worker quality is $q$; we make the simplifying assumption that this is independent of the employer type, since no signaling takes place in the pooling equilibrium.\footnote{Note this may not be completely accurate, since wage negotiation is endogenous, and employers' type may be partially revealed by wage negotiations.}    
Thus welfare in the pooling equilibrium is:
\begin{align}
  (\mu v_H + (1 - \mu) v_L) q.
\end{align}
Welfare in a separating equilibrium dominates welfare in a pooling equilibrium when 
\begin{equation}
\label{eq:welfarecond}
 \alpha_H q_H + \alpha_L q_L > q, 
\end{equation}
where:
\[ \alpha_H = \frac{\mu v_H}{\mu v_H + (1 - \mu) v_L},\ \ \alpha_L = \frac{(1 - \mu)v_L }{\mu v_H + (1 - \mu) v_L}. \]
We can interpret $\alpha_H$ (resp., $\alpha_L$) as the {\em value-weighted} fraction of high (resp., low) types in the marketplace.  
As one simple case where Equation~\ref{eq:welfarecond} holds, suppose that:
\[ \mu q_H + (1 - \mu) q_L \geq q. \]
In other words, this implies that the average quality of hired workers in the separating equilibrium is at least as high as the average quality of hired workers in the pooling equilibrium.  
For example, in supply-constrained markets, where most workers are hired, we would expect that this condition approximately holds.  
As long as $q_H > q_L$, then since $v_H > v_L$ it is straightforward to verify that Equation~\ref{eq:welfarecond} holds.

The preceding discussion provides some insight into how to interpret the suggested welfare increase in our setting.  
Essentially, the welfare gains to high type employers are sufficiently high to offset the welfare losses of low type employers, either because there is a sufficient differentiation in quality ($q_H \gg q_L$), or there is a sufficient mass of high type employers ($\alpha_H \gg \alpha_L$). 
